\subsection{L’OS et le Réseau : Les obstacles à la distribution}

Au-delà de l'optimisation algorithmique, l'obtention d'une latence faible et déterministe se heurte à un obstacle externe : l'imprévisibilité du système d'exploitation. En effet, l'ordonnanceur peut suspendre arbitrairement les \textit{threads} critiques du \textit{matching engine} au profit de tâches système. Ces changements de contexte (\textit{context switches}) imposent un cycle coûteux : la mise en pause de l'unité d'exécution, l'éviction des données courantes du cache au profit de celles du nouveau processus, puis, lors de la reprise, le rechargement nécessaire des caches CPU. Ce mécanisme de "va-et-vient" détruit la localité des données et accroît considérablement la gigue (\textit{jitter}) du programme.

Parallèlement au logiciel, le matériel intervient via les requêtes d'interruption (\textit{Interrupt ReQuests} ou \textit{IRQs}). Ces signaux forcent le processeur à suspendre instantanément le flux applicatif pour traiter des événements asynchrones, tels que l'arrivée de paquets réseau ou des accès disque, générant des pics de latence imprévisibles. La recherche d'une performance de l'ordre de la nanoseconde exige donc une stratégie d'isolation totale, visant à soustraire le moteur de trading à l'influence de l'ordonnanceur et des bruits système.

%Références:
%\begin{itemize}
%    \item Systems Performance, Brendan Gregg (2020)
%    \item  FIX Trading Community (2019). Simple Binary Encoding Specification
%\end{itemize}

\subsubsection{Configuration et exécution haute performance}

Afin de résoudre les context switches imposés par l'ordonnanceur, une technique consiste à isoler chaque \textit{thread} de l'application sur un cœur, c'est le \textit{CPU pinning}. En veillant à ce qu'un seul \textit{thread} soit présent par cœur, le CPU ne déclenche pas de changement de contexte. On veillera alors dans le chemin critique à ne jamais créer plus de \textit{threads} que de cœurs CPU. Imaginons un CPU possédant huit cœurs, le \textit{matching engine} possédera sept \textit{threads}, afin de garantir que les sept \textit{threads} critiques du moteur de \textit{matching} ne subissent aucune préemption, l'isolation des cœurs est configurée au niveau du chargeur d'amorçage via la directive \texttt{isolcpus} au sein du fichier \textit{/etc/default/grub}:

\begin{lstlisting}[caption={Configuration de l'isolation des cœœurs dans GRUB}]
GRUB_CMDLINE_LINUX_DEFAULT="isolcpus=1,2,3,4,5,6,7 nohz_full=1,2,3,4,5,6,7"
\end{lstlisting}

Note: Le huitième cœur est réservé à l'OS.

Une fois les ressources réservées, l'exécution du programme Java est réalisée en forçant l'affinité matérielle avec l'utilitaire \texttt{taskset}. Cette commande assure que la machine virtuelle Java (JVM) exploite exclusivement les segments de calcul isolés :

\begin{lstlisting}[caption={Exécution avec affinité CPU et optimisation NUMA}]
taskset -c 1-8 java -XX:+UseNUMA -cp . MonMatchingEngine
\end{lstlisting}

L'option \texttt{-XX:+UseNUMA} complète ce dispositif en optimisant l'allocation mémoire selon la topologie physique de la machine, minimisant ainsi les cycles d'attente lors des accès à la mémoire vive (RAM).

\subsubsection{Du Kernel Stack au Kernel Bypass: Vers le Zero-Copy réseau}

\paragraph{Les limites de la pile réseau standard}

Dans les systèmes de trading à haute fréquence, la traversée de la pile réseau représente la partie prépondérante de la latence de bout en bout, dépassant souvent le temps d'exécution de la logique métier. Comme l'illustre le tableau \ref{tab:serialization_latencies}, le délai de transmission pour un paquet standard passe de 120 000 ns en 100Base-T à seulement 1 200 ns en 10GBase-T. À ce niveau de performance, le coût logiciel devient prohibitif : la pile réseau standard du noyau Linux impose une latence de traversée de l'ordre de plusieurs microsecondes. Dès lors, pour un message de 256 octets, le temps passé par le CPU à copier les données entre les domaines mémoire (du \textit{Ring Buffer} de la NIC vers le \textit{Kernel Buffer}, puis vers le \textit{User Space}) excède mathématiquement le temps de traitement d'un ordre de bourse, souvent inférieur à 500 ns. 

\begin{table}[h!]
\centering
\begin{tabular}{@{}lrr@{}}
\toprule
\textbf{Bande Passante} & \textbf{Débit} & \textbf{Durée de transmission par paquet (ns)} \\ \midrule
100Base-T & 100 Mbps & 120 000 ns \\
1000Base-T (1G) & 1 Gbps & 12 000 ns \\
10GBase-T & 10 Gbps & 1 200 ns \\
100GBase-T & 100 Gbps & 120 ns \\
400GBase-T & 400 Gbps & 30 ns \\ \bottomrule
\end{tabular}
\caption{Durée de transmission par paquet selon la bande passante et le débit - Source Low Latency High-speed Systems - The Kernel Bypass Approach Ijaz Ahmad}
\label{tab:serialization_latencies}
\end{table}

Le transit d'un paquet du fil vers l'application subit une cascade d'indirections. À la réception, la carte réseau (NIC) génère une interruption matérielle (IRQ), forçant le processeur à suspendre le \textit{thread} applicatif pour exécuter le gestionnaire d'interruption du noyau. Ce mécanisme induit un double coût :

\begin{itemize}
    \item \textit{Cache Pollution} : le changement de contexte vers le mode noyau évince les lignes de cache dédiées à la logique métier, provoquant un besoin de les recharger à chaque retour au \textit{thread} métier
    \item \textit{Data Copying} : le paquet est copié du tampon de la NIC vers le \textit{kernel space}
\end{itemize}

\begin{figure}[H] % [H] force l'image à rester exactement ici
    \centering % Centre l'image sur la page
    \includegraphics[width=0.8\textwidth]{Kernel-Stack.png}
    \caption{Flux de réception des paquets via le Kernel - Source : Low Latency High-speed Systems - The Kernel Bypass Approach Ijaz Ahmad} % Légende sous l'image
    \label{fig:kernel-stack} % Label pour faire référence à l'image dans le texte
\end{figure}

\begin{figure}[H] % [H] force l'image à rester exactement ici
    \centering % Centre l'image sur la page
    \includegraphics[width=0.8\textwidth]{Kernel-Bypass.png}
    \caption{Représentation d'un flux de réception de paquets avec kernel bypass - Source : Low Latency High-speed Systems - The Kernel Bypass Approach Ijaz Ahmad} % Légende sous l'image
    \label{fig:kernel-bypass} % Label pour faire référence à l'image dans le texte
\end{figure}

\paragraph{La révolution du Kernel Bypass et le modèle Zero-Copy :}
Le \textit{Kernel Bypass} (via DPDK ou des piles propriétaires comme Solarflare Onload) court-circuite intégralement la couche logicielle de l'OS. La carte réseau écrit les données directement dans une zone mémoire partagée, accessible par l'application en mode utilisateur. L'impact de cette approche est triple :

\begin{itemize}
    \item Le thread dédié au réseau est épinglé sur un cœur CPU isolé et scrute en permanence les descripteurs de réception de la NIC. On élimine ainsi 100\% des interruptions et des changements de contexte.
    
    \item L'application gère elle-même les couches de transport (ex: pile TCP/UDP en espace utilisateur), ce qui permet de pré-allouer des buffers et d'éviter toute allocation dynamique sur le chemin critique.
    
    \item En alignant ces buffers sur les frontières des lignes de cache (64 octets), le \textit{matching engine} peut consommer le message dès sa réception sans aucun cycle d'horloge gâché par une copie intermédiaire.
\end{itemize}

Cette approche transforme l'interface I/O réseau en un simple accès mémoire, réduisant la gigue logicielle.

\subsubsection{Optimisation des messages: Simple Binary Encoding (SBE)}
Le passage d'un protocole textuel (FIX, JSON) ou binaire à structure variable (Google Protocol Buffers) vers SBE répond à une exigence de déterminisme à la nanoseconde.

\begin{itemize}
    \item \textbf{Élimination du coût de parsing et de l'allocation :} 
    La transformation de flux ASCII en objets Java (ex: \texttt{String} ou \texttt{Double}) est une hérésie en basse latence. Elle déclenche des allocations sur le \textit{heap}, augmentant la pression sur le \textit{Garbage Collector} et générant des pauses \textit{Stop-The-World}. SBE adopte une approche \textit{Allocation-Free} : les données sont lues directement depuis des \textit{DirectBuffers} (souvent mappés via \textit{Agrona} que nous présenterons ultérieuerement), sans instanciation d'objets intermédiaires. 
    
    \item \textbf{Philosophie \textit{Copy-Free} et accès séquentiel:} 
    Le design de SBE repose sur le principe de l'accès \textit{streaming}. Contrairement aux protocoles nécessitant un \textit{backtracking} ou des indirections pour accéder aux champs variables, SBE impose que les champs à longueur variable (comme les chaînes de caractères) soient placés en fin de message. Cette contrainte garantit que le processeur accède à la mémoire de manière strictement séquentielle. Cette linéarité est optimale pour le \textbf{L1 Data Prefetcher}. Le CPU anticipe le chargement des prochaines \textit{cache lines} (64 octets) dans le cache L1 avant même que le code ne les sollicite, évitant ainsi des cycles d'attente (\textit{stalls}) de pipeline.
\end{itemize}